\begin{proposition}\label{sga3:XXI:6.7.2}
Let \(\Delta\) be a system of simple roots, and put
\[
 E_\Delta(\cR)=
 \{u\in\operatorname{Aut}(\cR)\mid u(\Delta)=\Delta\}.
\]
Then \(\operatorname{Aut}(\cR)\) is the semidirect product of
\(W(\cR)\) by \(E_\Delta(\cR)\).
\end{proposition}

\begin{proof}
This follows at once from the fact that \(W(\cR)\) acts simply
transitively on the systems of simple roots and from the fact that,
if \(\Delta\) is a system of simple roots of \(\cR\), then
\(u(\Delta)\) is a system of simple roots for every automorphism
\(u\) of \(\cR\).
\end{proof}

We shall later give a simpler description of \(E_\Delta(\cR)\) in
the case of reduced irreducible root data.\footnote{Editors' note:
When \(\cR\) is simply connected or adjoint, see \SGARefLink{sga3:XXI:7.4.5}{7.4.5}.}

\begin{definition}\label{sga3:XXI:6.7.3}
We denote by \(\operatorname{Aut}^{s}(\cR)\)\footnote{Editors' note:
The superscript \(s\) is intended to suggest ``semisimple.''}
the set of \(u\in\operatorname{Aut}(\cR)\) for which the following
diagram is commutative:
\begin{center}
\begin{tikzcd}[column sep=huge,row sep=large]
\cR \arrow[rr,"u"] \arrow[dr]
  && \cR \arrow[dl] \\
& \operatorname{rad}(\cR) &
\end{tikzcd}
\end{center}
We write
\[
 E_\Delta^{s}(\cR)
 =E_\Delta(\cR)\cap\operatorname{Aut}^{s}(\cR).
\]
\end{definition}

\begin{remark}\label{sga3:XXI:6.7.4}
If \(u\in\operatorname{Aut}(\cR)\), then
\(u\in\operatorname{Aut}^{s}(\cR)\) if and only if
\[
 (u-\operatorname{id})(M)\subset\VR(R).
\]
In particular,
\[
 W(\cR)\subset\operatorname{Aut}^{s}(\cR).
\]
\end{remark}

\begin{proposition}\label{sga3:XXI:6.7.5}
For every system of simple roots \(\Delta\), the group
\(\operatorname{Aut}^{s}(\cR)\) is the semidirect product of
\(W(\cR)\) by \(E_\Delta^{s}(\cR)\).
\end{proposition}

By functoriality, every automorphism of \(\cR\) induces an
automorphism of \(\operatorname{ad}(\cR)\). Thus there is a
canonical homomorphism
\[
 \operatorname{Aut}(\cR)\longrightarrow
 \operatorname{Aut}(\operatorname{ad}(\cR)).
\]

\begin{lemma}\label{sga3:XXI:6.7.6}
The homomorphism
\[
 \operatorname{Aut}^{s}(\cR)\longrightarrow
 \operatorname{Aut}(\operatorname{ad}(\cR))
\]
is injective.
\end{lemma}

\begin{proof}
Let \(u\) be an automorphism of \(M\) such that
\[
 (u-\operatorname{id})(M)\subset\VR(R)
\]
and
\[
 {}^tu(\alpha^\ast)=\alpha^\ast
 \qquad(\alpha^\ast\in R^\ast).
\]
For every \(x\in M\), one has
\[
 (\alpha^\ast,u(x)-x)
 =({}^tu(\alpha^\ast)-\alpha^\ast,x)=0.
\]
Hence \(u(x)-x=0\), by \ref{sga3:XXI:1.2.5}.
\end{proof}

\begin{lemma}\label{sga3:XXI:6.7.7}
The group \(\operatorname{Aut}^{s}(\cR)\) is finite.
\end{lemma}

\begin{proof}
It suffices to prove that \(\operatorname{Aut}(\cR)\) is finite when
\(\cR\) is adjoint. In that case \(M\) is generated by \(R\), so
every automorphism of \(\cR\) is determined by the permutation it
induces on \(R\).
\end{proof}
